\subsection{完备空间}

在一致空间 X 中，Cauchy 滤子未必有极限点。

例。1）考虑序列 \((u_{n})\)（在有理直线 \(\mathbf{Q}\) 上），它由下式定义：

\[
\text { by } \quad u _ {n} = \sum_ {p = 0} 2 ^ {- p (p + 1) / 2}. \quad \text { If } \quad m > n \quad \text { we   have }\tag{1}
\]

\[
\left| u _ {m} - u _ {n} \right| \leqslant 2 ^ {- n (n + 3) / 2}
\]

因而 \((u_{n})\) 是 Cauchy 序列。但这一序列在 \(\mathbf{Q}\) 中没有极限；因为若有理数 \(a / b\) 是 \((u_{n})\) 的极限，则由 (1)，对所有 \(n\) 我们应有

\[
\left| a / b - h _ {n} / 2 ^ {n (n + 1) / 2} \right| \leqslant 1 / 2 ^ {n (n + 3) / 2}
\]

其中 \(h_n\) 是（依赖于 \(n\) 的）整数；即

\[
\left| a. 2 ^ {n (n + 1) / 2} - b h _ {n} \right| \leqslant b. 2 ^ {- n}
\]

对所有 n 成立。现在，对一切 n，这一不等式的左端都是整数，因而只要 n 大于某个整数 \(n_{0}\)（满足 \(b < 2^{n_{0}}\)），它就必须为零；于是应有 \(a/b = u_{n}\)（对所有 \(n > n_{0}\)），这是荒谬的。

\begin{enumerate}
\def\labelenumi{\arabic{enumi})}
\setcounter{enumi}{1}
\tightlist
\item
  设 X 是无穷集合，并考虑 X 上有限分拆的一致结构（§ 2 第 2 目注 1）。X 上的每个超滤子 F 都是
\end{enumerate}

关于这一一致结构的 Cauchy 滤子。因为若 \((\mathbf{A}_{i})\) 是 X 的一个有限分拆，且

\[
\mathbf {V} = \bigcup_ {i} (\mathbf {A} _ {i} \times \mathbf {A} _ {i})
\]

是其对应的一致邻域，则至少有一个 \(A_{i}\) 属于 F（第一章 § 6 第 4 目命题 5 的推论），而 \(A_{i}\) 是 V-小集。另一方面，X 是无穷离散空间，因而不是紧的，从而 X 上存在不收敛的超滤子。

定义 3 完备空间是指每个 Cauchy 滤子都收敛的一致空间。

因此，在完备空间中，每个 Cauchy 序列（第 1 目）都收敛。例 在离散一致空间 X 上，Cauchy 滤子是平凡超滤子（第一章 § 6 第 4 目），因而收敛；从而 X 是完备的。

由第 1 目的定义 2 与定义 3 以及第 1 目的命题 2，我们推出下述命题，它称为 Cauchy 判别准则：

命题 6 设 \(\mathfrak{F}\) 是集合 X 上的滤子，\(f\) 是 X 到完备一致空间 \(\mathbf{X}'\) 的映射。则 \(f\) 关于 \(\mathfrak{F}\) 有极限的充要条件是：\(\mathfrak{F}\) 在 \(f\) 下的像是 Cauchy 滤子基。

这一判别准则表明了完备空间在涉及极限概念的一切问题中的重要性：如果一个函数在完备空间中取值，我们就可以在事先不知道极限值的情况下证明极限的存在；倘若极限的定义是我们掌握的仅有的收敛判别准则，这就不可能做到。

细于完备空间的一致结构的一致结构，未必是完备空间的一致结构（习题 2）。然而，我们有下述命题：

命题 7 设 \(\mathcal{U}_1, \mathcal{U}_2\) 是集合 \(X\) 上的两个一致结构，\(\mathfrak{C}_1, \mathfrak{C}_2\) 分别是这些一致结构诱导的拓扑。假设 \(\mathcal{U}_1\) 细于 \(\mathcal{U}_2\)，且存在 \(\mathcal{U}_1\) 的一个一致邻域基本系，其中的元素作为 \(X \times X\) 的子集在拓扑 \(\mathfrak{C}_2 \times \mathfrak{C}_2\) 中是闭的。则滤子 \(\mathfrak{F}\)（\(X\) 上的）在拓扑 \(\mathfrak{C}_1\) 中收敛的充要条件是：它在一致结构 \(\mathcal{U}_1\) 中是 Cauchy 滤子，且在拓扑 \(\mathfrak{C}_2\) 中收敛。

这些条件显然是必要的，因为 \(\mathfrak{G}_2\) 粗于 \(\mathfrak{G}_1\)。反之，设条件满足，并设 \(x\) 是 \(\mathfrak{F}\) 关于 \(\mathfrak{G}_2\) 的极限点；我们将证明 \(x\) 是 \(\mathfrak{F}\) 关于 \(\mathfrak{G}_1\) 的极限。设 \(V\) 是 \(\mathfrak{U}_1\) 的一个对称一致邻域，它在拓扑 \(\mathfrak{G}_2 \times \mathfrak{G}_2\) 中是闭的。由假设，\(\mathfrak{F}\) 含有一个集合 \(M\)，它是 \(V\)-小集；于是若 \(x' \in M\)，我们就有 \(M \subset V(x')\)。但 \(V(x')\) 在拓扑 \(\mathfrak{G}_{2}\) 中是闭的；因此，\(x\) 位于 \(\mathbf{M}\) 关于 \(\mathfrak{G}_{2}\) 的闭包中，故必属于 \(\mathbf{V}(x')\)。由此得出 \(\mathbf{M} \subset \mathbf{V}(x)\)，命题得证。

推论 在命题 7 的条件下，若 \(\mathfrak{U}_2\) 是完备空间的一致结构，则 \(\mathfrak{U}_1\) 也是完备空间的一致结构。

因为此时每个关于 \(U_{1}\) 的 Cauchy 滤子都是关于 \(U_{2}\) 的 Cauchy 滤子，从而在拓扑 \(G_{2}\) 中收敛。

注意，当 \(\mathcal{C}_1 = \mathcal{C}_2\) 时（§ 1 第 2 目命题 2 的推论 2），命题 7 的推论的假设成立。

